\section{Disjoint Set Union}

\begin{frame}
\frametitle{Disjoint Set Union}
\begin{block}{Problem: Disjoint Set Union}
    Are vertices $u$ and $v$ in the same connected component?
\end{block}

\begin{center}
\begin{tikzpicture}
[  
    level 1/.style={sibling distance=25mm},
    level 2/.style={sibling distance=15mm},
    edge from parent/.style={thick}
]
    % Left tree (red)
    \begin{scope}[edge from parent/.style={draw=red}]
        \node [draw, circle, red] {1}
            child {node [draw, circle, red] {3}}
            child {
                node [draw, circle, red] {2}
                child {node [draw, circle, red] {4}}
            };
    \end{scope}
    
    % Right tree (orange)
    \begin{scope}[edge from parent/.style={draw=orange}]
        \node [draw, circle, orange] at (6,0) {5}
            child {
                node [draw, circle, orange] {7}
                child {node [draw, circle, orange] {6}}
                child {node [draw, circle, orange] {8}}
            };
    \end{scope}
\end{tikzpicture}
\end{center}
\end{frame}

\begin{frame}[fragile]
\frametitle{Disjoint Set Union}

We can just DFS!

\begin{lstlisting}[language=C++, basicstyle=\small]
int a, b;
bool reach_b = false;
void dfs(int u, int p) {
    if(u == b) {
        reach_b = true;
    }
    for(auto v: adj[u]) {
        if(v == p) {
            continue;
        }
        dfs(v, u);
    }
}
// ...
dfs(a, 0);
\end{lstlisting}

Time Complexity: $\mathcal{O}(V + E)$
\end{frame}

\begin{frame}
\frametitle{Disjoint Set Union}

It quite slow... How about changing the adjacency list to an array.

\begin{block}{Definition: Parent Array}
    parent[u] := direct parent of vertex $u$ but if $u$ has no parent parent[u] will store itself.
\end{block}

\begin{center}
\begin{tikzpicture}
[  
    level 1/.style={sibling distance=25mm},
    level 2/.style={sibling distance=15mm},
    edge from parent/.style={thick}
]
    \begin{scope}[edge from parent/.style={draw=red}]
        \node [draw, circle, red] {1}
            child {node [draw, circle, red] {3}}
            child {
                node [draw, circle, red] {2}
                child {node [draw, circle, red] {4}}
            };
    \end{scope}
    
    \begin{scope}[edge from parent/.style={draw=orange}]
        \node [draw, circle, orange] at (6,0) {5}
            child {
                node [draw, circle, orange] {7}
                child {node [draw, circle, orange] {6}}
                child {node [draw, circle, orange] {8}}
            };
    \end{scope}
\end{tikzpicture}
\end{center}

\begin{center}
\begin{tabular}{|c||c|c|c|c|c|c|c|c|}
    \hline
    index & 1 & 2 & 3 & 4 & 5 & 6 & 7 & 8\\
    \hline
    \hline
    parent & 1 & 1 & 1 & 2 & 5 & 7 & 5 & 7\\
    \hline
\end{tabular}
\end{center}

\end{frame}

\begin{frame}[fragile]
\frametitle{Disjoint Set Union}
\begin{center}
\begin{tabular}{|c||c|c|c|c|c|c|c|c|}
    \hline
    index & 1 & 2 & 3 & 4 & 5 & 6 & 7 & 8\\
    \hline
    \hline
    parent & 1 & 1 & 1 & 2 & 5 & 7 & 5 & 7\\
    \hline
\end{tabular}
\end{center}

So, we can easily write \verb|find_root| function to find the root of component.

\begin{lstlisting}[language=C++, basicstyle=\small]
int parent[N];
int find_root(int u) {
    if(u == parent[u]) {
        return u;
    }
    return find_root(parent[u]);
}
\end{lstlisting}

Time Complexity: $\mathcal{O}(V)$

\end{frame}

\begin{frame}[fragile]
\frametitle{Disjoint Set Union}
\begin{center}
\begin{tabular}{|c||c|c|c|c|c|c|c|c|}
    \hline
    index & 1 & 2 & 3 & 4 & 5 & 6 & 7 & 8\\
    \hline
    \hline
    parent & 1 & 1 & 1 & 2 & 5 & 7 & 5 & 7\\
    \hline
\end{tabular}
\end{center}

Or we can just use while loop!

\begin{lstlisting}[language=C++, basicstyle=\small]
int parent[N];
int find_root(int u) {
    while(u != parent[u]) {
        u = parent[u];
    }
    return u;
}
\end{lstlisting}

Time Complexity: $\mathcal{O}(V)$

\end{frame}

\begin{frame}[fragile]
\frametitle{Disjoint Set Union}
It will be very slow if height of the tree is large.

\begin{center}
\begin{verbatim}
    1 - 2 - 3 - 4 - 5 - 6 - 7 - 8 - 9 - 10
\end{verbatim}
\end{center}

By calling \verb|find_root(10)| is \textcolor{red}{very slow}! So, we have to try to optimize this!

\end{frame}

